\documentclass{article}
\usepackage{../../style/header}

\begin{document}

Field theory is very hard, Galois theory is the gymnastics of recovering information about fields from the much easier subjects of linear algebra, and finite groups acting on finite sets.

\begin{definition}
    The \textbf{degree} of a field extension $K\subset L$ is $[L:K] := \dim_K L$ as a $K$-vector space. A field extension is called finite if its degree is.
\end{definition}

From now on, unless otherwise said, all field extensions are finite.

\begin{proposition}[Tower law]
    If $K\subset F \subset L$ is a tower of extensions, then $[L:K]=[L:F][F:K]$.\begin{proof}
        If $B = \{b_1,\ldots,b_n\}$ and $C=\{c_1,\ldots,c_m\}$ is a $K$-basis for $F$ and an $F$-basis for $L$ respectively, then $\{b_ic_j \mid 1\leq i \leq n, \ 1\leq j \leq m\}$ is a $K$-basis for $L$.
    \end{proof}
\end{proposition}

All field homomorphisms $K\rightarrow L$ are injective, as the kernel is a proper ideal of $K$. So, instead of $\Hom(K,L)$, I'll write $\Emb(K,L)$ for the set of field homomorphisms $K\rightarrow L$. By the first isomorphism theorem, a field \textit{endomorphism} $K\rightarrow K$ must be an isomorphism, so I'll write $\Aut(K)$ for $\Emb(K,K)$.

\begin{definition}
    If $K\subset F,L$ are field extensions, then the set of \textbf{$K$-embeddings} from $F$ to $L$ is \[
    \Emb_K(F,L) = \{f\in \Emb(F,L) \mid f\vert_K = \id_K\}
    \]
\end{definition}

In Galois theory, $K$ is almost always the prime subfield, and so this distinction often doesn't matter. When there are multiple embeddings $K\rightarrow F,L$, by writing $K\subset F,L$ I am choosing a canonical embedding, and specifying some $f\in \Emb_K(F,L)$ is asking for this map to commute with the canonical embeddings, in which case things like the cardinality of this set may depend on my choice of canonical embedding.

I will avoid this complication by almost always having a fixed \textit{ground field} $k$ of which every field $K$ is an extension, so I am really working in the category of finite field extensions $k\subset K$. Here $k$ is the special initial object, and I will try to avoid using $k$ as an object in propositions when this property isn't necessary.

\begin{proposition}[A2]
    If $K \subset L$ is a field extension, then there is a \textit{right} action of $\Aut_k(K)$ on $\Emb_k(K,L)$ by composition, and this action is free.
    \begin{proof}
        As all $f\in \Emb_k(K,L)$ are injective, if there exists $\sigma \in \Aut_k(K)$ with $f\sigma = f$, then $\sigma=\id_K$.
    \end{proof}
\end{proposition}

\subsection{Minimal polynomials}

\begin{definition}
    If $K\subset L$ is a field extension with $a\in L$, then the \textbf{minimal polynomial} $f$ of $a$ is the monic generator of the kernel of the evaluation map $\ev_a:K[X]\rightarrow L$ by $\ev_a(\phi) = \phi(a)$. Write $K(a) = \im\ev_a$, which is clearly the unqiue intermediate extension generated by $a$.
\end{definition}

As $f(a)=0$, we have $K(a)=K[a]$, giving us an explicit inverse and a $K$-basis.

\begin{lemma}
    With the above setup, there is a canonical bijection $\Emb_K(K(a),L) = \{\text{roots of $f$ in $L$}\}$.
    \begin{proof}
        Such an embedding $\phi\in \Emb_K(K(a),L)$ is fully determined by $\phi(a)$, for which the possibilities are exactly the roots of $f$ in $L$, use the isomorphism theorem to get the exact embedding.
    \end{proof}
\end{lemma}

\begin{proposition}
    Suppose $f\in K[X]$ with $K\subset L_1$ is generated by the roots of $f$ and $f$ splits completely in $K\subset L_2$, then $\Emb_K(L_1,L_2)\neq\emptyset$.
    \begin{proof}
        Do induction on $[L_1:K]$, where the base case is trivial. Now, choose some $a\in L_1\setminus K$ with $f(a)=0$ and let $g\in K[X]$ be its minimal polynomial, then $g\mid f$ splits completely in $L_2$ and we may find an embedding $K(a)\subset L$, so it suffices to show $\Emb_{K(a)}(L_1,L_2)\neq\emptyset$, which is the inductive hypothesis and $[L_1:K(a)]<[L_1:K]$.
    \end{proof}
\end{proposition}

\begin{corollary}
    So if $f\in K[X]$, then repeatedly adjoining roots of $f$ to $K$ gives a field $K\subset L$ which is unique up to non-unique $k$-isomorphism.
\end{corollary}

\subsection{Categorics}

\begin{proposition}[A4]
    For all finite extensions $K\subset F,L$ there exists a diagram: \[\begin{tikzcd}
        K \arrow[r, phantom, "\subset"] \arrow[d, phantom, sloped, "\subset"]  & F \arrow[d, phantom, sloped, "\subset"] \\
        L \arrow[r, phantom, "\subset"] & \Omega
    \end{tikzcd}\]
\begin{proof}
    Go by induction on $[F:K]$, the base case has $F=K$ where taking $\Omega = L$ works. 
    
    Now suppose $[F:K]>1$, there will exist some $a\in F\setminus K$, which will have minimal polynomial $f\in K[X]$, let $L'$ be the splitting field of $f\in L[X]$, then \[\Emb_k(k(a),L') = \{\text{roots of $f$ in $L'$}\}\neq \emptyset\] and we can form the daigram:
    \[\begin{tikzcd}
        K \arrow[r, phantom, "\subset"] \arrow[d, phantom, sloped, "\subset"] & K(a) \arrow[r, phantom, "\subset"] \arrow[d, phantom, sloped, "\subset"] & F \\
        L \arrow[r, phantom, "\subset"] & L'
    \end{tikzcd}\]
    where, as $K\neq K(a)$, the tower law says $[F:K(a)] < [F:K]$.
\end{proof}
\end{proposition}

This looks like a pushout, but because we are asking for canonical embeddings $F,L\subset \Omega$ there is no universal property.

It's an exercise on one of the problem sheets to show fibered products, which are given by intersections, and ``framed'' pushouts, which are pushouts in the category of subfields of some $\Omega$, and given by the field-theoretic product $FL$, exist. This is (A3)

\begin{definition}
    $k\subset K$ is a \textbf{normal extension} if for all $K\subset\Omega$ the action $\Emb_k(K,\Omega)\curvearrowleft \Aut_k(K)$ is transitive, i.e.~every $x:K\hookrightarrow \Omega$ is just the canonical inclusion of $\sigma(K)\subset K$ into $\Omega$.
\end{definition}

An extension is thus normal iff for all $x_1,x_2\in \Emb_k(K,\Omega)$, $x_1(K)\subseteq x_2(X)$.

\begin{lemma}
    If $k\subset K$ is normal, and $k\subset F \subset K$ is an intermediate field, then $F\subset K$ is normal.
    \begin{proof}
        Given $x,y\in\Emb_F(K,\Omega)\subseteq\Emb_k(K,\Omega)$ there exists some $\sigma\in\Aut_k(K)$ such that $x=y\sigma$, this $\sigma$ must be in $\Aut_F(K)$.
    \end{proof}
\end{lemma}

\begin{proposition}
    Let $k\subset K$ be normal and $k\subset F \subset K \subset \Omega$ be a tower, then the natural restriction: \[
    \Emb_k(K,\Omega)\rightarrow \Emb_k(F,\Omega)
    \] is surjective, i.e.~any embedding $x:F\rightarrow\Omega$ extends to some $\tilde{x}:K\rightarrow\Omega$.\begin{proof}
    For any $x\in\Emb_k(F,\Omega)$, by (A4) I can form the following diagram: \[
    \begin{tikzcd}
        & \Omega \arrow[r, "y"] & \Omega' \\
        k \arrow[r, phantom, sloped, "\subset"] & F\arrow[r, phantom, sloped, "\subset"] \arrow[u, "x"] & K \arrow[u, "z", swap] \arrow[ul, dashed, "\tilde{x}", swap]
    \end{tikzcd} 
    \] but as $k\subset K$ is normal, I must be able to write $z:K\rightarrow\Omega'$ as $y\circ\sigma$ for some $\sigma\in\Aut_k(K)$. This $\sigma$ (composed with the inclusion $K\subset\Omega$) is exactly the required $\tilde{x}$.
\end{proof}
\end{proposition}

From this, it is immediately obvious that the natural injective composition with $K\subset\Omega$: \[
    \Emb_k(F,K)\rightarrow\Emb_k(F,\Omega)
\] is also surjective.

\begin{proposition}[A1]
    For all extensions $k\subset K,L$ the set $\Emb_k(K,L)$ has cardinality at most $[K:k]$.
    \begin{proof}
        Once again going by induction on $[K:k]$ where the base case is trivial.

        Supposing $[K:k]>1$, pick $\alpha\in K\setminus k$ and consider the restriction map: \[
            \rho:\Emb_k(K,L)\rightarrow \Emb_k(k(\alpha),L)
        \] Now for any $x\in \Emb_k(k(\alpha),L)$, the pre-image $\rho^{-1}(x)=\Emb_x(K,L)$ has cardinality $\leq[K:k(\alpha)]$, so I can consider the sum: \[
            \abs{\Emb_k(K,L)} = \sum_{x\in\Emb_k(k(\alpha),L)}\abs{\rho^{-1}(x)} \leq [k(\alpha):k][K:k(\alpha)] = [K:k]\qedhere
        \]
    \end{proof}
\end{proposition}

\begin{lemma}
    For an extension $k\subset K$ and $H\leq \Aut(K/k)$ a finite subgroup, $[K:K^H]\leq \abs{H}$.
    \begin{proof}
        If $G = \{\sigma_1,\ldots,\sigma_n\}$, we need to show all sets $\{a_1,\ldots,a_{n+1}\}\subset L$ are linearly dependent over $K^H$. Let $\{a_1,\ldots,a_k\}$ consider the vectors: \[
        \mathbf{a}_1=\begin{pmatrix}
            \sigma_1(a_1) \\
            \vdots \\
            \sigma_n(a_1)
        \end{pmatrix},
        \ldots,
        \mathbf{a}_{n+1}=\begin{pmatrix}
            \sigma_1(a_{n+1}) \\
            \vdots \\
            \sigma_n(a_{n+1})
        \end{pmatrix}
        \] and wlog let $\{\mathbf{a}_1,\ldots,\mathbf{a}_k\}$ be the smallest nontrivial linearly dependent set, so we have $x_i\neq 0$ such that $x_1\mathbf{a}_1 + \cdots + x_k\mathbf{a}_k = 0$, wlog take $x_1=1$. Applying some $\sigma\in \Aut(K/k)$ will permute the rows of this linear system, but will always keep $x_1$ the same, this must always be the same system, otherwise we could subtract and find one of strictly smaller size. Thus all the $x_i\in K^H$ and equation corresponding to $\sigma_i=e$ is as required.
    \end{proof}
\end{lemma}

\begin{corollary}[A5]
    Under the same assumptions, $[K:K^H]=\abs{H}$, and the obvious inclusion $H\subset\Aut(K/K^H)$ is an isomorphism.\begin{proof}
        (A1) already says $\abs{\Aut(K/K^H)}\leq [K:K^H]$, by the previous lemma we know $[K:K^H]\leq \abs{H}$, and by assumption we know $\abs{H}\leq \abs{\Aut(K/K^H)}$.
    \end{proof}
\end{corollary}

\begin{definition}
    The extension $k\subset K$ is \textbf{separable} if for all towers of subfield: \[
    k \subset K_1 \subsetneq K_2 \subset K
    \] there exists some field $\Omega$ with at least two distince $K_1$-embeddings $x,y:K_2\rightarrow\Omega$.
\end{definition}

\begin{lemma}
    If $k\subset K$ is separable, and $k\subset F \subset K$ is a tower, then $k\subset F$ and $F\subset K$ are separable.
    \begin{proof}
        Any pair of intermediate field of $k\subset F$ or $F\subset K$ are also intermediate field of $k\subset K$.
    \end{proof}
\end{lemma}

\begin{theorem}[Fundamental theorem of Galois theory]
    If $k\subset K$ is Galois (normal and separable) then the there is a bijective, order-reversing map between field and subgroups as follows: \[
        H\leq G \longmapsto K^H \qquad k\subset F \subset K \longmapsto \Aut(K/F)
    \]\begin{proof}
        We know the images are fields and subgroups respectively, and the fact that these maps reverse the inclusion ordering is obvious. $H = \Aut(K/K^H)$ is exactly (A5). As all $k\subset F\subset K$ are also Galois, we just have to prove $k':=K^{\Gal(K/k)} = k$ we will go by showing that for all $k\subset \Omega$: $\Emb_k(k',\Omega)$ is just the inclusion $k'\subset K \subset \Omega$, where separability now implies the final result. Consider the tower: \[
            k\subset k' \subset K \subset \Omega
        \] we know any $\Emb_k(k',\Omega)$ is the restriction of some $\Emb_k(K,\Omega)$ but as $k\subset K$ is normal, this corresponds to composing the inclusion $K\subset \Omega$ with some $\sigma\in \Gal(K/k)$, by construction this is just $k'\subset K\subset \Omega$.
    \end{proof}
\end{theorem}

We now want to translate the categorical definitions of normal and separable extensions into field theoretic statements.

\subsection{Normal extensions}

\begin{theorem}
    For a finite extension $k\subset K$ tfae: \begin{enumerate}
        \item For all irreducible $f\in k[X]$, either $f$ has no roots in $K$ or $f$ splits completely in $K$;\footnote{This is a standard definition of a normal extension}
        \item There exists some $f\in k[X]$ such that $K$ is the splitting field of $f$ over $k$;
        \item $k\subset K$ is normal.
    \end{enumerate}
    \begin{proof}
        $(1\Rightarrow2)$ If $K=k(a_1,\ldots,a_n)$ and $f_i$ is the minimal polynomial of $a_i$ over $k$ respectively, $K$ is clearly the splitting field of $f_1\cdots f_n$,

        $(2\Rightarrow3)$ If $k\subset K$ is the splitting field of $f$ with root $a\in K$, then for any $K\subset \Omega$ and $x:K\rightarrow \Omega$, $x(a)$ will always be a root of $f$ in $\Omega$ and as $K$ is generated by all the roots of $f$, $x(K)\subset K \subset \Omega$.

        $(3\Rightarrow 1)$ Now if $k\subset K$ is normal, for any irreducible $f\in k[X]$ let $K\subset\Omega$ be the splitting field of $f$ over $K$. If $f$ has any root $a\in K$ then we can form the tower $k\subset k(a)\subset K\subset \Omega$ giving the bijection: \[
            \{\text{roots of $f$ in $\Omega$}\} = \Emb_k(k(a),\Omega) = \Emb_k(k(a),K) = \{\text{roots of $f$ in $K$}\}\qedhere
        \] 
    \end{proof}
\end{theorem}

\begin{lemma}
    For all $k\subset K$ there exists some $K\subset L$ such that $k\subset L$ is normal.\begin{proof}
        Let $K=k(a_1,\ldots,a_n)$, with $f_i$ the minimal polynomial of each $a_i$ over $k$ respectively, now let $L$ be the splitting field of $f_1\cdots f_n$ over $K$, and thus also over $k$.
    \end{proof}
\end{lemma}

\subsection{Separable extensions}

\begin{definition}
    For $k\subset K$ define the \textbf{separable degree}\[
        [K:k]_s := \abs{\Emb_k(K,\Omega)}
    \] where $k\subset K \subset \Omega$ is normal extension that exists by the previous lemma.
\end{definition}

\begin{lemma}
    The separable degree $[K:k]_s$ does not depend on the choice of $\Omega$.\begin{proof}
        Let $\Omega_1,\Omega_2$ be such normal extensions of $k$, then by (A4) let $\Omega'$ be a field containing both of them, we know any embedding $K\rightarrow \Omega'$ must send $K$ to the normal subfield containg it, i.e. \[
            \Emb_k(K,\Omega_1) = \Emb_k(K,\Omega) = \Emb_k(K,\Omega_2)\qedhere
        \]
    \end{proof}
\end{lemma}
Separable extensions are now exactly those $k\subset K$ such that for all towers $k\subset K_1\subset K_2\subset K$, if $[K_1:K_2]_s=1$, then $K_1=K_2$.

\begin{proposition}
    For a tower $k\subset K\subset L$ we have \[
    [L:k]_s=[L:K]_s[K:k]_s
    \]
    \begin{proof}
        Take $L\subset \Omega$ such that $k\subset \Omega$ is normal, we know the restriction: \[
        \Emb_k(\Omega,\Omega)\rightarrow\Emb_k(L,\Omega)\xrightarrow{\rho} \Emb_k(K,\Omega)
        \] is surjective, so we can count\[
        [L:k]_s := \abs{\Emb_k(L,\Omega)} = \sum_{x\in\Emb_k(K,\Omega)}\abs{\rho^{-1}(x)} = [K:k]_s[L:K]_s
        \] as the previous lemma says the size of $\rho^{-1}(x) = \Emb_K(L,\Omega)$ doesn't depend on the choice $x:K\rightarrow \Omega$.
    \end{proof}
\end{proposition}

\begin{definition}
    A polynomial $f\in K[X]$ is \textbf{separable} if it has $n=\deg(f)$ distinct roots in its splitting field. An element $a\in K$ is then called separable if its minimal polynomial is.
\end{definition}

\begin{proposition}
    An irreducible polynomial $f\in K[X]$ is inseparable iff $K$ has characteristic $p>0$ and there exists some irreducible $h\in K[X]$ with $f(X)=h(X^p)$.
\end{proposition}

\begin{lemma}
    If $K\subset K(a)$ is separable, then $a$ is separable over $K$.
    \begin{proof}
        Let $f\in K[X]$ be the minimal polynomial of $a$ over $K$, and suppose it isn't separable, then we have $f(X) = h(X^p)$ for some irreeucible $h$, and $b=a^p$ is a root of $h$, so we can form the tower $K\subset K_1 = K(b) \subset K_2 = K(a)$ where, by the tower law \[
        p \deg(h) = \deg(f) = [K_2:K] = [K_2:K_1][K_1:K] = [K_2:K_1]\deg(h)
        \] so $[K_2:K_1] = p$ and $(X-a)^p = X^p-b\in K_1[X]$ is the minimal polynomial of $a$ over $K_1$, so $[K_2:K_1]_s = 1$, and thus $K\subset K(a)$ is inseparable.
    \end{proof}
\end{lemma}

\begin{theorem}
    A finite extension $k\subset K$ is separable iff $[K:k]_s = [K:k]$.
    \begin{proof}
        ($\impliedby$) If $[K:k]_s = [K:k]$, then by applying two tower laws to \[
        k \subset K_1 \subset K_2 \subset K
        \] and by (A1), $[-:-]_s\leq [-:-]$, if $[K_2:K_1]_s = 1$, so if $[K_2:K_1]_s = 1$, we get \[
        [K:K_2][K_2:K_1][K_1:k] = [K:K_2]_s[K_1:k]_s \leq [K:K_2][K_1:k]
        \] so $[K_2:K_1]=1$ and $K_1 = K_2$.

        ($\implies$) Now, if $k\subset K$ is separable, pick $a\in K\setminus k$ and consider the tower $k\subset k(a) \subset K$, we have $[k:k(a)]$ separable, so by the previous lemma we know $a$ is separable over $k$, so \[
        [k(a):k] = \deg(f) = \abs{\{\text{roots of }f\}} = [k(a):k]_s
        \] and as $[K:k(a)]<[K:k]$, we may inductively assume $[K:k(a)]_s = [K:k(a)]$, and the tower law finished the proof.
    \end{proof}
\end{theorem}

\begin{corollary}
    The following list of statements immediately hold: \begin{enumerate}
        \item For all towers $k\subset K \subset L$, is $k\subset K$ and $K\subset L$ are separable, then $k\subset L$ is separable;
        \item $k\subset K$ is separable iff every $a\in K$ is separable over $k$;
        \item if $a$ is separable over $k$, then $k\subset k(a)$ is separable;
        \item if $k\subset K$ is the splitting field of a separable polynomial, then $k\subset K$ is separable.
    \end{enumerate}
\end{corollary}

\subsection{Discriminants}

\begin{definition}
    If $K$ is a field and $f\in K[X]$ is monic, then in the splitting field we can write \[
    f(X) = (X-a_1)(X-a_2)\cdots(X-a_n)
    \] then when $S_n$ acts on these roots, the \textbf{discriminant} of $f$ is the invariant quantity \[
    \Delta_f = \prod_{i<j}(a_i-a_j)^2
    \] so we must be able to write this in term of elementary symmetric polynomials, for small $n$ this looks like \begin{gather*}
        f(X) = X^2 - e_1X + e_2 \quad \implies \quad \Delta_f = e_1^2 - 4e_2 \\
        f(X) = X^3 - e_1X^2 + e_2X - e_2 \quad \implies \quad \Delta_f = e_1^2e_2^2 - 4e_1^3e_2 - 4e_2^3 + 18e_1e_2e_3 - 27e_3^2 \\
        \text{so } f(X) = X^3 + pX + q \quad \implies \quad \Delta_f = -4p^3 - 27q^2
    \end{gather*}
\end{definition}

\begin{proposition}
    If $K$ is a field of characteristic $\neq 2$, and $f(X)\in K[X]$ is a monic separable polynomial with splitting field $K\subset L$, then $\Gal(L/K)\subset A_n$ iff $\Delta_f$ is a square in $K$.\begin{proof}
        We just use the fact that the square root of $\Delta_f$ is in $K$ iff it is fixed by $\Gal(L/K)$, and is defined so to be fixed by $A_n$, but not $S_n$.
    \end{proof}
\end{proposition}

\subsection{Primitive element theorem}

This is Stacks project lemma 9.19.1, and very much not categorical.

\begin{proposition}
        For all finite field extensions $K\subset L$, tfae: \begin{enumerate}
        \item $L$ is \textbf{simple} over $K$, i.e.~there exists some $\alpha \in L$ such that $L=K(\alpha)$;
        \item there are finitely many intermediary subfields $K\subset F \subset L$.
    \end{enumerate}
    \begin{proof}
        ($\Rightarrow$) suppose we have a subfield $K\subset F \subset K(\alpha)$ with $f$ the minimal polynomial of $\alpha$ over $F$, then $\deg(f)=[K(\alpha):F]$ and if $f=x^d + a_{d-1}x^{d-1}+\cdots + a_0$, then $K(a_0,\ldots,a_{d-1})\subset F$ and they have the same degree in $K(\alpha)$, so are equal. All such $f$ must be one of the finitely many factors of the minimal polynomial of $\alpha$ in $K$ when viewed in $K(\alpha)$.

        ($\Leftarrow$) If we are working with finite fields then we can always find primitive elements as generators of the cyclic groups $L^\times$. Now assume $K$ is infinite, and has only finitely many subfields $K_1,\ldots, K_n$, these are all proper subspaces so we can find $\alpha \not\in K_i$ for all $i$, this is our primitive element.
    \end{proof}
\end{proposition}

\begin{proposition}
        If $K\subset L$ is separable, then the above conditions hold.
    \begin{proof}
        Let $k\subset L\subset\Omega$ be a normal extension---so $K\subset L \subset \Omega$ is as well---and consider an enumeration of the finite set $\Emb_K(L,\Omega) = \{\sigma_1,\ldots,\sigma_n\}$, if $K\subset L$ is separable, then $[L:K]=n$ and for all $i\neq j$ we have $V_{ij} = \ker(\sigma_i-\sigma_j)\subsetneq L$ from which we can find a primitive element as in the previous proposition.
    \end{proof}
\end{proposition}

When actually trying the find primitive elements, it is probably easier to write out how the Galois group acts on a nice basis, and find an combination that isn't fixed.

\subsection{Biquadratic extensions}

\subsection{Finite fields}

\begin{proposition}
    For $p>0$ prime and $m>0$ an integer, there exists a field $\bF_q$, with $q=p^m$, unique up to isomorphism, with $\Gal(\bF_q/\bF_p)\cong C_m$ generated by $\Fr_p:a\mapsto a^p$.\begin{proof}
        Suppose such a field $F$ exists, then the $\abs{F^\times}=q-1$, so every element is a solution of $X^q-X\in\bF_p[X]$, so $F$ is the splitting field, unique up to isomorphism, of $X^q-X$. By the Jacobian criterion this is separable.

        Now, if $F$ is the splitting field of $X^q-X\in\bF_p[X]$, we already know this is separable, so $F$ has at least $q$ elements. The clever part now is to observe that the solutions of $X^q-X$ infact form a field, so there are at most $q$ elements and we are done.

        Now I can consider $\Fr_p:\bF_q\rightarrow \bF_q$, this has degree dividing $m$ in $\abs{\Gal(\bF_q/\bF/p)}=m$, if it had lower degree, $p^{m'}$, then all of $\bF_q$ would be a solution to a separable degree $p^{m'}$ polynomial.
    \end{proof}
\end{proposition}

You can do some estimation and analysis to find, with the inclusion-exclusion counting for the Galois correspondence for $\bF_q$, that there is always an irreducible monic polynomial of degree $n$ over $\bF_p$.

\begin{corollary}
    Every extension of finite field is normal and separable with cyclic Galois group.\begin{proof}
        We know the size of a finite extension $\bF_p\subset \bF_q$, so all extensions are of the form $\bF_{p^r}\subset \bF_{p^{rm}}$, but as $\bF_p\subset \bF_{p^{rm}}$ is normal and separable, so is $\bF_{p^r}\subset \bF_{p^{rm}}$.
    \end{proof}
\end{corollary}

If $P$ is a, assumed commutative, monoid and $K$ a field, then a function $\chi:P\rightarrow K$ is a (multiplicative) character if it is a monoid homomrophism into $K$ with multiplication.

\begin{lemma}[Dedekind independence theorem]
    A set of $n$ pairwise distinct nontrivial characters $\{\chi_1,\ldots, \chi_n\}$ from $P$ to $K$ is linearly independent in the $K$-vector space of functions $P\rightarrow K$.\begin{proof}
        Do induction on $n$, the base case $n=1$ is trivial. Assume, for a contradiction, we have a linear relation $\lambda_1\chi_1 + \cdots \lambda_n\chi_n = 0$, then fix $p\in P$ with $\chi_1(p)\neq \chi_2(p)$. By the multiplicativity of the $\chi_i$s, we get another linear relation $\lambda_1\chi_1(p)\chi_1 + \cdots + \lambda_n\chi_n(p)\chi_n = 0$, if we combine these to cancel out the $\chi_1$ coefficient, then the $\chi_2$ coefficient is definitely nonzero, but by induction it must be zero, a contradiction, hence the $\chi_i$s are linearly independent.
    \end{proof}
\end{lemma}

\begin{theorem}
    Let $f(X)\in\bZ[X]$ be degree $n$ monic; $\bQ\subset K$ a splitting field; $p$ a prime such that the reduction $f_p\in\bF_p[X]$ of $f$ mod $p$ has $n$ distinct roots; $\bF_p\subset F$ a splitting field of $f_p$. Let $\lambda_1,\ldots,\lambda_n\in K$ be the roots of $f$, and $R = \bZ[\lambda_1,\ldots,\lambda_n]\subset K$, then we have \begin{enumerate}
        \item there exists a ring homomorphism $\psi:R\rightarrow F$;
        \item any such homomorphism maps $\lambda_1,\ldots,\lambda_n$ bijectively to the roots of $f_p$;
        \item the action $\Hom_{\mathbf{Ring}}(R,F)\curvearrowleft\Gal(K/\bQ)$ is transitive.
    \end{enumerate}
    \begin{proof}
        First, $R$ is a free finitely generated $\bZ$-module, generated by the monomials $\lambda_1^{k_1}\lambda_2^{k_2}\cdots\lambda_n^{k_n}$ where $0\leq k_i\leq n-1$ for all $i$, we can always use $f$ to reduce a higher exponent.

        Now, as $\bZ[\frac{1}{p}]$ isn't a finitely generated $\bZ$-module, it cannot be a $\bZ$-submodule of $R$. So, we can find a maximal ideal $(p)\subset \fm \subset R$ and the field $\bF_p\subset R/\fm$ will be generated by the roots of $f_p$, so isomorphic to the splitting field $F$.

        If $u_1,\ldots,u_d$ is a basis of $R$ as a $\bZ$-module, then they are also $\bQ$-linearly independent (clear denominators), and from the problem sheet, we know $\bQ\subset \bQ R\subset K$ must be a field, as it contains all the roots of $f$, it must be all of $K$, hence $u_1,\ldots,u_d$ is a $\bQ$-basis of $K$.

        Finally, as $\abs{\Gal(K/\bQ)} = [K:\bQ] = d$, write $\Gal(K/\bQ) = \{\sigma_1,\ldots,\sigma_d\}$, then for all ring homomorphisms $\psi':R\rightarrow F$, if $\psi'$ is not one of $\psi\sigma_1,\ldots,\psi\sigma_d$, then by considering these are characters $R\rightarrow F$, for all $1\leq j\leq r$, we can find $\lambda_i\in F$ not all zero, such that $(\lambda_0\psi' + \lambda_1\psi\sigma_1 + \cdots \lambda_d\psi\sigma_d)(u_j) = 0$ (we have $d+1$ unknowns to solve a system of $d$ equations), which contradicts Dedekind independence.
    \end{proof}
\end{theorem}

In particular, there exists an element $\sigma\in \Gal(K/\bQ)$ such that $\Fr_p\psi = \psi\sigma$, we call such an element a \textbf{Frobenius lift}.

If $f$ is separable, then $\gcd(f,f')$ is $1$, and so I can find $\phi,\psi\in\bZ[X]$ such that $f\phi + f'\psi = m\in \bZ$, when $p\nmid m$, and I take this equation modulo $p$, I get the reduction of $f$ and $f'$ are coprime, so $f$ mod $p$ is separable for all but finitely many $p$.

\begin{corollary}
    If $f_p$ factors as a product of irreducible factors of degree $n_1,\ldots, n_k$, then $\Gal(K/\bQ)$ contains an element of cycle shape $(n_1)(n_2)\cdots (n_k)$.
    \begin{proof}
        The fixed points of $\Fr_p$ are $\bF_p$, and so $\Fr_p$ will be a nontrivial permutation of the roots of each irreducible factor of $f_p$, thus a Frobenius lift will be a permutation of this cycle shape.
    \end{proof}
\end{corollary}

If we accept a few lemmata from group theory.

\begin{lemma}
    If $G\leq S_n$ is transitive, $G$ contains a transposition, and an $(n-1)$-cycle, then $G=S_n$.
\end{lemma}

\begin{lemma}
    If $p$ is a prime, $G\leq S_p$, and $G$ contains a transposition and a $p$-cycle, then $G=S_p$.
\end{lemma}

\subsection{Cyclotomic polynomails}

\begin{lemma}
    For all integers $n>0$, the cyclotomic polynomial $\Phi_n(X)$ is in $\bZ[X]$.
    \begin{proof}
        We have $\Phi_n(X)\in\bQ(\mu_n)[X]$, as $\bQ\subset \bQ(\mu_n)$ has Galois group a subgroup of $\Aut\mu_n = (\bZ/n\bZ)^\times$, the exponents of the primitive roots of unity, which fixes $\Phi_n(X)$, which must be in $\bQ[X]$. We can write \[X^n-1 = \prod_{d\mid n}\Phi_d(X)\] and now we are done by the Gauss lemma and induction on $n$.
    \end{proof}
\end{lemma}

\begin{theorem}
    For all $n>0$, the canonical injection $\Gal(\bQ(\mu_n)/\bQ)\hookrightarrow \Aut\mu_n$ is an isomorphism.
    \begin{proof}
        If $p\nmid n$ is a prime, then $X^n-1\in\bF_p[X]$ is separable with splitting field $F$, so as in the previous theorem, we have a group isomorphism $\psi:\mu_n(\bZ[\zeta_n])\xrightarrow{\sim} \mu_n(F)$. If we let $\sigma\in \Gal(\bQ(\mu_n)/\bQ)$ be a Frobenius lift, then this acts as $z\mapsto z^p$, and the set of these for all such $p$ generates $\Aut\mu_n$.
    \end{proof}
\end{theorem}

\begin{corollary}
    $\Phi_n(X)$ is irreducible, and $[\bQ(\mu_n):\bQ] = \varphi(n)$.
\end{corollary}

\subsection{Kummer theory}

\begin{theorem}
    For $n>0$ an integer, and $K$ a field of characteristic $p\nmid n$ containg all $n$th roots of unity, a field extensions $K\subset L$ is of the form $L=K(\beta)$ for some nonzero $\beta^n=a\in K$ iff $K\subset L$ is normal and separable with cyclic Galois group whose order divides $n$
    \begin{proof}
        ($\Rightarrow$) Choose $\beta$ with $\beta^n=a$, then consider the function $\Gal(L/K)\rightarrow \mu_n$ by $\sigma\mapsto \sigma(\beta)/\beta$, that this is well-defined, and in fact an injective group homomorphism are all direct.

        ($\Leftarrow$) Take $x\in L$, $\sigma\in \Gal(L/K)=C_k$, and $\zeta\in\mu_n$ both generators, and consider \[
            \alpha = x + \zeta^{-1}\sigma(x) + \zeta^{-2}\sigma^2(x) + \cdots + \zeta^{1-k}\sigma^{k-1}(x)
        \] By Dedekind independce, there exists some $x\in L$ such that $\alpha$ is nonzero, and satisfies \[
        \sigma(\alpha) = \sigma(x) + \sigma^{-1}\sigma^2(x) + \cdots + \sigma^{1-k}\sigma^k(x) = \zeta(x + \zeta^{-1}\sigma(x) + \cdots + \zeta^{1-k}\sigma^{k-1}(x)) = \zeta\alpha
        \] So if $a = \alpha^n$, then $a\in K$ as it's fixed by $\Gal(L/K)$, and as $\Gal(L/K)$ acts on $K(\alpha)$ we have \[
        [L:K] = \abs{G} \leq [K(\alpha):K] \leq [K:L]\qedhere
        \]
    \end{proof}
\end{theorem}
And so if we assume $K$ has enough roots of unity, by the Galois correspondence, a solvabe $\Gal(L/K)$, with composition series \[
1 = G_0 \lhd G_1 \lhd \cdots \lhd G_r = \Gal(L/K)
\] will give a tower of subfields \[
L = L_0 \supset L_1 \supset \cdots \supset L_r = K
\] where each $L_{i+1}\subset L_i$ is normal and has Galois group $\Gal(L_{i+1}/L_i) \cong G_{i+1}/G_i$ cyclic, and so we can write $K\subset L$ as a composition of adjoining $n_i$th roots. This relationship is iff.

\begin{corollary}
    For $n\geq 5$ the polynomial $f(X) = (X-X_1)\cdots (X-X_n)\in K(X_1,\ldots,X_n)[X]$ is not solvable by radicals.
    \begin{proof}
        If we write $\sigma_1,\ldots,\sigma_n$ for the elementary symmetric polynomials, then $K(\sigma_1,\ldots,\sigma_n)$ will be the splitting field of $f$, from algebra 3 we know $K(\sigma_1,\ldots,\sigma_n) = K(X_1,\ldots,X_n)^{S_n}$, and so our extension has Galois group $S_n$, which is not solvable.
    \end{proof}
\end{corollary}


\end{document}